\section{선별 문제 풀이 I}
\label{sec:finite-field-solutions-a}

\subsection*{문제 \thechapter.2: $\F_4$의 연산표}

\begin{solution}
$\alpha^2+\alpha+1=0$이고 표수가 $2$이므로
\[
  \alpha^2=\alpha+1,
  \qquad
  \alpha^3=\alpha(\alpha+1)=\alpha^2+\alpha=1.
\]
따라서
\[
  \F_4=\{0,1,\alpha,\alpha+1\}.
\]

덧셈표는 다음과 같다.
\[
\begin{array}{c|cccc}
+&0&1&\alpha&\alpha+1\\
\hline
0&0&1&\alpha&\alpha+1\\
1&1&0&\alpha+1&\alpha\\
\alpha&\alpha&\alpha+1&0&1\\
\alpha+1&\alpha+1&\alpha&1&0
\end{array}
\]

곱셈표는 다음과 같다.
\[
\begin{array}{c|cccc}
\cdot&0&1&\alpha&\alpha+1\\
\hline
0&0&0&0&0\\
1&0&1&\alpha&\alpha+1\\
\alpha&0&\alpha&\alpha+1&1\\
\alpha+1&0&\alpha+1&1&\alpha
\end{array}
\]

$0^4=0$, $1^4=1$이고 $\alpha^3=1$이므로 $\alpha^4=\alpha$이다. 또한 $\alpha+1=\alpha^2$이므로
\[
  (\alpha+1)^4=\alpha^8=\alpha^2=\alpha+1.
\]
따라서 네 원소가 모두 $X^4-X$의 근이다.
\end{solution}

\subsection*{문제 \thechapter.3: $\F_8$의 거듭제곱 계산}

\begin{solution}
관계식 $\alpha^3=\alpha+1$을 반복하여 사용한다.
\[
\begin{array}{c|c}
 k&\alpha^k\\
\hline
0&1\\
1&\alpha\\
2&\alpha^2\\
3&\alpha+1\\
4&\alpha^2+\alpha\\
5&\alpha^2+\alpha+1\\
6&\alpha^2+1\\
7&1
\end{array}
\]

$\F_8^\times$의 위수는 $7$이고 $\alpha\ne1$이다. 위수 $7$인 군에서 항등원이 아닌 원소의 위수는 $7$이므로 $\alpha$는 생성원이다.

마지막으로
\[
  \alpha^2+\alpha+1=\alpha^5.
\]
따라서 역원은 지수의 합이 $7$이 되도록 택하여
\[
  (\alpha^2+\alpha+1)^{-1}
  =(\alpha^5)^{-1}=\alpha^2
\]
이다.
\end{solution}

\subsection*{문제 \thechapter.5: $X^4-X$의 분해}

\begin{solution}
$\F_2$의 두 원소 $0,1$이 근이므로
\[
  X(X+1)\mid X^4-X.
\]
나눗셈을 하거나 차수가 $2$ 이하인 일계수 기약다항식을 분류하면 남은 인수는 $X^2+X+1$이다. 따라서
\[
  X^4-X=X(X+1)(X^2+X+1).
\]
표수 $2$에서는 $-X=X$이므로 직접 전개해도
\[
  X(X+1)(X^2+X+1)=X^4+X=X^4-X
\]
이다.

$X$와 $X+1$의 근은 $\F_2$의 원소이고, $X^2+X+1$의 두 근은 $\F_4\setminus\F_2$의 원소이다. 따라서 세 인수의 근 전체가 정확히 $\F_4$이다.
\end{solution}

\subsection*{문제 \thechapter.7: 교집합과 합성체}

\begin{solution}
최대공약수와 최소공배수 공식을 적용하면
\[
  \F_{p^6}\cap\F_{p^{10}}
  =\F_{p^{\gcd(6,10)}}=\F_{p^2}
\]
이고
\[
  \F_{p^6}\F_{p^{10}}
  =\F_{p^{\operatorname{lcm}(6,10)}}=\F_{p^{30}}.
\]
따라서
\[
  [\F_{p^{30}}:\F_{p^6}]=5,
  \qquad
  [\F_{p^{30}}:\F_{p^{10}}]=3.
\]
교집합 위의 차수는
\[
  [\F_{p^6}:\F_{p^2}]=3,
  \qquad
  [\F_{p^{10}}:\F_{p^2}]=5
\]
이다.
\end{solution}

\subsection*{문제 \thechapter.9: $\F_8$의 프로베니우스 궤도}

\begin{solution}
상대 프로베니우스는 제곱사상이다. 먼저
\[
  \alpha\longmapsto\alpha^2
  \longmapsto\alpha^4=\alpha^2+\alpha
  \longmapsto\alpha^8=\alpha
\]
이므로 궤도는
\[
  \{\alpha,\alpha^2,\alpha^4\}
\]
이다. 따라서 최소다항식은
\[
  (X-\alpha)(X-\alpha^2)(X-\alpha^4)
  =X^3+X+1.
\]

이제 $\beta=\alpha+1=\alpha^3$라 하자. 지수를 법 $7$로 계산하면
\[
  \beta^2=\alpha^6=\alpha^2+1,
  \qquad
  \beta^4=\alpha^{12}=\alpha^5=\alpha^2+\alpha+1.
\]
다음 제곱은 다시 $\beta$이므로 궤도 길이는 $3$이다. 또한
\[
  \beta^3=\alpha^2,
  \qquad
  \beta^2=\alpha^2+1
\]
이어서
\[
  \beta^3+\beta^2+1=0.
\]
따라서 $\beta$의 최소다항식은
\[
  X^3+X^2+1
\]
이다.
\end{solution}

\subsection*{문제 \thechapter.10: 기약다항식의 개수}

\begin{solution}
모비우스 반전 공식을 차례로 적용한다.
\[
  N_2(1)=2.
\]
또한
\[
  N_2(2)=\frac{2^2-2}{2}=1,
  \qquad
  N_2(3)=\frac{2^3-2}{3}=2.
\]
$n=4$에서는 $\mu(4)=0$이므로
\[
  N_2(4)=\frac{2^4-2^2}{4}=3.
\]
마지막으로
\[
  N_3(2)=\frac{3^2-3}{2}=3.
\]
\end{solution}

\subsection*{문제 \thechapter.12: $X^q-X$의 근들이 이루는 체}

\begin{solution}
$f=X^q-X$라 하자. $q=p^n$이고 표수가 $p$이므로
\[
  f'=qX^{q-1}-1=-1.
\]
따라서 $f$는 중근을 갖지 않고 대수적 폐포에서 정확히 $q$개의 서로 다른 근을 갖는다.

그 근 집합을 $S$라 하자. $a,b\in S$이면 반복된 프로베니우스 공식으로
\[
  (a+b)^q=a^q+b^q=a+b,
\]
\[
  (-a)^q=-a,
  \qquad
  (ab)^q=a^qb^q=ab.
\]
$b\ne0$이면
\[
  (b^{-1})^q=(b^q)^{-1}=b^{-1}.
\]
따라서 $S$는 덧셈, 덧셈역원, 곱셈과 영이 아닌 원소의 역원에 닫혀 있으며 체이다. 근이 정확히 $q$개이므로 $|S|=q$이다.
\end{solution}

\subsection*{문제 \thechapter.14: 유한체의 포함관계}

\begin{solution}
먼저 $\F_{p^m}\subseteq\F_{p^n}$이라 하자. 탑 법칙으로
\[
  n=[\F_{p^n}:\F_p]
  =[\F_{p^n}:\F_{p^m}]m
\]
이므로 $m\mid n$이다.

반대로 $n=rm$이라 하자. $a\in\F_{p^m}$이면 $a^{p^m}=a$이다. 이 등식을 반복하면
\[
  a^{p^n}=a^{p^{rm}}=a.
\]
따라서 $a$는 $X^{p^n}-X$의 근이고 $\F_{p^n}$에 속한다. 모든 $a$에 대해 성립하므로
\[
  \F_{p^m}\subseteq\F_{p^n}.
\]
\end{solution}
